%% \newcommand{\MyContourText}[2][\huge]{\contour{black}{#1 \protect \textcolor{white}{#2}}}

\newcommand{\MyLectureTitlePage}[0]{
    \setbeamertemplate{background canvas}{\includegraphics[height=\paperwidth,width=\paperheight,angle=-90]{motherboard}}
    \begin{frame}
      \titlepage
    \end{frame}
}

\newcommand{\MyLectureTOCOld}[0]{\begin{frame}[shrink]%% [allowframebreaks]
    \frametitle{Table of contents}
    \begin{multicols}{2}
      \tableofcontents
    \end{multicols}
  \end{frame}
}

\newcommand{\MyLectureTOC}[0]{
\begin{frame}[shrink]%% [allowframebreaks]
    \frametitle{Table of contents}
    \begin{multicols}{2}
      \tableofcontents[hideallsubsections]
    \end{multicols}
  \end{frame}
}

\makeatletter
\def\equalsfill{$\m@th\mathord=\mkern-7mu
\cleaders\hbox{$\!\mathord=\!$}\hfill
\mkern-7mu\mathord=$}
\makeatother

\newcommand{\MyMiscellanea}[2]{
  \section{Miscellanea}
  \subsection{Corrections and additions since last lecture.}
  \begin{frame}{}
    \begin{columns}
      \begin{column}{0.5\textwidth}
        \begin{itemize}
          #1
        \end{itemize}
      \end{column}
      \begin{column}{0.5\textwidth}
        \includegraphics[width=\textwidth]{misc_socialmedia_klein}
      \end{column}
    \end{columns}
    \vspace{\baselineskip}
    #2
  \end{frame}
}

\newcommand{\MyComputerOrganizationBook}[0]{
  \section{The Book}
\subsection{Class text}
  \begin{frame}{Computer Organization, Brief Edition  \cite{sterlingBrief}}
    \begin{columns}
      \begin{column}{0.5\textwidth}
        \begin{itemize}
        \item Specified in the syllabus
        \item Based on extracts from other texts \cite{null2010essentials,computerSystemsWarford}
\item Original texts are available on-line, from the publisher, from book sellers
        \end{itemize}
Occassional references in Sterling's to figures and chapters in the orginial texts.
      \end{column}
      \begin{column}{0.5\textwidth}
        \includegraphics[width=\textwidth]{computerOrganization}
      \end{column}
    \end{columns}
    \vspace{\baselineskip}
  \end{frame}
\subsection{Original texts}
  \begin{frame}{The Essentials Computer Organization Architecture}
    \begin{columns}
      \begin{column}{0.5\textwidth}
        Some sections are complete copies of the original texts.
        \linebreak
        ~
        \linebreak
        \MyChapterXref{2}{3}{Boolean Algebra and Digital Logic}
        \cite{null2010essentials}
      \end{column}
      \begin{column}{0.5\textwidth}
        \includegraphics[width=\textwidth]{computerOrganizationNull}
      \end{column}
    \end{columns}
    \vspace{\baselineskip}
  \end{frame}
  \begin{frame}{Computer Systems}
    \begin{columns}
      \begin{column}{0.5\textwidth}
        Some sections are complete copies of the original texts.
        \footnotesize
        \begin{itemize}
        \item \MyChapterXref{1}{3}{Information Representation}
        \item \MyChapterXref{4}{4}{Computer Architecture}
        \item \MyChapterXref{5}{5}{Computer Architecture}
        \item \MyChapterXref{6}{6}{Compiling to the Assembly Level}
        \end{itemize}
        \normalsize
        Source of HOL6, LG1, and other notations \cite{computerSystemsWarford}
      \end{column}
      \begin{column}{0.5\textwidth}
        \includegraphics[width=\textwidth]{computerSystemsWarford}
      \end{column}
    \end{columns}
    \vspace{\baselineskip}
  \end{frame}

}


%% \newcommand{\MyFrontMatterPresentation}[0]{\MyFrontMatter[courseImage]}

\newcommand{\MyChapterXref}[3]{
  Sterling Chap. #1 from original Chap. #2: #3
}

\newcommand{\MyLogicalOperators}[0]{
\begin{frame}{Logical operators}
  \begin{columns}
    \begin{column}{0.5\textwidth}
      Simple in concept.  Complex in operation.
      \begin{itemize}
      \item AND (\&) only 1 when both are 1
      \item OR ($\vert$) 1 when either are 1
      \item XOR (\textasciicircum) 1 when only one is 1
      \item NOT (\textasciitilde) flips 1s to 0s, and 0s to 1s
      \end{itemize}
    \end{column}
    \begin{column}{0.5\textwidth}
      \begin{center}
        \includegraphics[width=\textwidth]{bitops01-01}
      \end{center}
    \end{column}
  \end{columns}\vspace{\baselineskip}
A 0 can be FALSE, and \textasciicircum FALSE can be TRUE.
\end{frame}
}

\newcommand{\MyExam}[2]{
  \section{Exam \##1}
  \begin{frame}{You will have #2.}
  \begin{columns}
    \begin{column}{0.5\textwidth}
      \begin{itemize}
      \item Be sure and put your name on the exam.
\item Your work must be your own.
\item When finished bring your exam up to the front.
\item If you have a question, come up and ask.  I will answer for the entire class.
      \end{itemize}
    \end{column}
    \begin{column}{0.5\textwidth}
      \begin{center}
        \includegraphics[width=0.8\textwidth]{examCheatSheet}
        \linebreak
        An example of a 3D cheat sheet.
      \end{center}
    \end{column}
  \end{columns}\vspace{\baselineskip}
 We'll have a quick review after the break.
\end{frame}
\MyBreakTime
\section{Review}
\begin{frame}{}
  \begin{columns}
    \begin{column}{0.5\textwidth}
      A quick review of the exam before we start today's lecture.
    \end{column}
    \begin{column}{0.5\textwidth}
      \begin{center}
        \includegraphics[width=0.8\textwidth]{3x9}
        \MyImageCredit{blog:xkcdExams}
      \end{center}
    \end{column}
  \end{columns}\vspace{\baselineskip}
\end{frame}
}


\newcommand{\MyLectureDateOne}[0]{\MyMondayA}
\newcommand{\MyLectureDateTwo}[0]{\MyMondayB}
\newcommand{\MyLectureDateThree}[0]{\MyMondayC}
\newcommand{\MyLectureDateFour}[0]{\MyMondayD}
\newcommand{\MyLectureDateFive}[0]{\MyMondayE}
\newcommand{\MyLectureDateSix}[0]{\MyMondayF}
\newcommand{\MyLectureDateSeven}[0]{\MyMondayG}
\newcommand{\MyLectureDateEight}[0]{\MyMondayH}
\newcommand{\MyLectureDateNine}[0]{\MyMondayI}
\newcommand{\MyLectureDateTen}[0]{\MyMondayJ}
\newcommand{\MyLectureDateEleven}[0]{\MyMondayK}
\newcommand{\MyLectureDateTwelve}[0]{\MyMondayL}
\newcommand{\MyLectureDateThirteen}[0]{\MyMondayM}
\newcommand{\MyLectureDateFourteen}[0]{\MyMondayN}
\newcommand{\MyLectureDateFifteen}[0]{\MyMondayO}
\newcommand{\MyLectureDateSixteen}[0]{\MyMondayP}
\newcommand{\MyLectureDateSeventeen}[0]{\MyMondayError}

%% \newcommand{\MyLectureTitleOne}[0]{Sections 1.1 and 1.3, Theory of bits}
%% \newcommand{\MyLectureTitleTwo}[0]{Sections 1.4 and 1.5, Theory of bits}
%% \newcommand{\MyLectureTitleThree}[0]{Sections 1.6, Theory of bits}
%% \newcommand{\MyLectureTitleFour}[0]{Test, Sections 2.1 and 2.2, Combinational circuits}
%% \newcommand{\MyLectureTitleFive}[0]{Sections 2.3 and 2.4, Combinational circuits}
%% \newcommand{\MyLectureTitleSix}[0]{Test, Section 3.1, Sequential circuits}
%% \newcommand{\MyLectureTitleSeven}[0]{Sections 3.2 and 3.3, Sequential circuits}
%% \newcommand{\MyLectureTitleEight}[0]{Test, Section 4.1 and 4.2, Computer architecture}
%% \newcommand{\MyLectureTitleNine}[0]{Sections 4.3 and 4.4, Computer architecture}
%% \newcommand{\MyLectureTitleTen}[0]{Section 5.1 and 5.2, Assembly language}
%% \newcommand{\MyLectureTitleEleven}[0]{Test, Section 5.3 and 5.4, Assembly language}
%% \newcommand{\MyLectureTitleTwelve}[0]{Test, Section 6.1 and 6.2, Assemblers and high order languages}
%% \newcommand{\MyLectureTitleThirteen}[0]{Sections 6.3 and 6.4, Assemblers and high order languages}
%% \newcommand{\MyLectureTitleFourteen}[0]{Sections 6.5, Assemblers and high order languages}
%% \newcommand{\MyLectureTitleFifteen}[0]{Test, Course review}
%% \newcommand{\MyLectureTitleSixteen}[0]{Final}
%% \newcommand{\MyLectureTitleSeventeen}[0]{\ThisShouldCauseTheCompiliationToBreak}

\newcommand{\MyLectureTitleOne}[0]{Sections 1.1 through 1.3, Information Representation}
\newcommand{\MyLectureTitleTwo}[0]{Sections 1.4 and 1.5, Information Representation}
\newcommand{\MyLectureTitleThree}[0]{Sections 1.6, Information Representation}
\newcommand{\MyLectureTitleFour}[0]{Test, and Section 2.1 Combinational Circuits}
\newcommand{\MyLectureTitleFive}[0]{Sections 2.2 through 2.4, Combinational Circuits}
\newcommand{\MyLectureTitleSix}[0]{Test, Section 3.1, Sequential Circuits}
\newcommand{\MyLectureTitleSeven}[0]{Sections 3.2 and 3.3, Sequential Circuits}
\newcommand{\MyLectureTitleEight}[0]{Test, Section 4.1 and 4.2, Computer Architecture}
\newcommand{\MyLectureTitleNine}[0]{Sections 4.3 and 4.4, Computer Architecture}
\newcommand{\MyLectureTitleTen}[0]{Review, Section 5.1 and 5.2, Assembly Language}
\newcommand{\MyLectureTitleEleven}[0]{Test, Section 5.3 and 5.4, Assembly Language}
\newcommand{\MyLectureTitleTwelve}[0]{Test, Section 6.1 and 6.2, Compiling to the Assembler Level}
\newcommand{\MyLectureTitleThirteen}[0]{Sections 6.3 and 6.4, Compiling to the Assembler Level}
\newcommand{\MyLectureTitleFourteen}[0]{Sections 6.5, Compiling to the Assembler Level}
\newcommand{\MyLectureTitleFifteen}[0]{Test, Course review}
\newcommand{\MyLectureTitleSixteen}[0]{Final}
\newcommand{\MyLectureTitleSeventeen}[0]{\ThisShouldCauseTheCompiliationToBreak}

\makeatletter
\newcommand*{\textoverline}[1]{$\overline{\hbox{#1}}\m@th$}
\makeatother

\newcommand{\MyChapter}[0]{Not defined}

\newcommand{\MyTitle}[3]{\title{CSC-223: Data Structures and Analysis of Algorithms\\Lecture \##1\\#2}
  \author{\MyContourText{Dr. Chuck Cartledge}}
  \date{\MyContourText{#3}}}

\newcommand{\MyPepAssemblerPage}[0]{195 (figure labeled:The Pep/9 instruction set at Level ISA3)}
\newcommand{\MyPepMnemonicPage}[0]{240 (figure labeled:The Pep/9 instruction set at Level Asmb5)}
\newcommand{\MyASCIIPage}[0]{31}

\newwrite\tempfile
\immediate\openout\tempfile=external-lists.tex

\newcounter{MyCounter}
\newcommand\ShowMyCounter{\stepcounter{MyCounter}\theMyCounter}

\newcommand{\MyOutcomeOne}[0]{Sets, Relations, and Functions}
\newcommand{\MyOutcomeTwo}[0]{Proof Techniques}
\newcommand{\MyOutcomeThree}[0]{Basics of Counting}
\newcommand{\MyOutcomeFour}[0]{Basic logic}
\newcommand{\MyOutcomeFive}[0]{Graph \& Trees}
\newcommand{\MyOutcomeSix}[0]{Discrete Probability}
\newcommand{\MyOutcomeSeven}[0]{Recurrence Relations}
\newcommand{\MyOutcomeEight}[0]{Boolean Algebra \& Expressions }
\newcommand{\MyOutcomeNine}[0]{Combinatorial Circuits}
